% arara: lualatex: { shell: yes, options: "-interaction nonstopmode -halt-on-error -file-line-error" }
% arara: biber
% arara: makeglossaries
% arara: lualatex: { shell: yes, options: "-interaction nonstopmode -halt-on-error -file-line-error" }
% arara: lualatex: { shell: yes, options: "-interaction nonstopmode -halt-on-error -file-line-error" }

% mgr.tex - root top
% version: 0.1
% TODO:
% - pierwsza strona - podmienic obrazki
% - zyciorys?
% - flagi draftu
% - impnattypo

\documentclass{eitidypl}

% %%%%%%%%%%%%%%%%%%%%%%%%%%%%%%%%%%%%%%%%%%%%%%%%%%%%%%%%%%%%%%%%%%%%%%%%%%%%%

\rokegz{2017 r.}
\rokak{2016/2017}
\semestr{Lato}
\stopien{II}
\kierunek{Telekomunikacja}
\instytut{telekomunikacji}
\jednostka{Instytut Telekomunikacji}
\specjalnosc{Telekomunikacja}
\typ{magisterska}
\autor{Jan Kowalski}
\nralbumu{123456}
\adresa{ul. Puławska 5/12}
\adresb{02-111 Warszawa}
\foto{foto.png}
\dataurodzenia{01 grudnia 1234r.}
\datarozpoczecia{1. stycznia 2000r.}
\opiekun{dr hab. inż. ...}
\tytul{Tytul twojej pracy}
\tytulen{English title}

\zyciorys{Urodziłem się 1 stycznia 1981r. w Warszawie...\\[3mm]
abc...\\[3mm]
def...\\[3mm]
Poniżej podpis:}

\streszczenie{Polskie streszczenie pracy...\\[3mm]
Dalsza część streszczenia...\\[3mm]
I coś jeszcze}

\streszczenieen{English abstract...\\[3mm]
Something more...\\[3mm]
And something else...}

\slowakluczowe{{słowa} {kluczowe} {pracy}}
\slowakluczoween{english, keywords}

%  %%%%%%%%%%%%%%%%%%%%%%%%%%%%%%%%%%%%%%%%%%%%%%%%%%%%%%%%%%%%%%%%%%%%%%%%%%%%

\usepackage{refcount}
%fonts
%\usepackage{times}
%\usepackage{lmodern} %lm
%computer modern ?
%http://tex.stackexchange.com/questions/24994/switching-back-to-computer-modern-within-xelatex

%caption – Customising captions in floating environments
%http://www.peteryu.ca/tutorials/publishing/latex_captions#set_caption_label_numbering_style
\ifthenelse{\boolean{strict}}{
\usepackage[format=hang,labelformat=simple,labelsep=colon,justification=raggedright,{font=footnotesize,sf},singlelinecheck=off]{caption}
}{
\usepackage[format=hang,labelformat=simple,labelsep=colon,justification=raggedright,{font=footnotesize},singlelinecheck=off]{caption}

}
\renewcommand{\thefigure}{\arabic{chapter}.\arabic{figure}}
\renewcommand{\thetable}{\arabic{chapter}.\arabic{table}}

\usepackage{float}
%mozna rozluznic above, below
\usepackage[section]{placeins}
%Set space between lines - should be default
%\usepackage{setspace}

%%%%%%%%%%%%%%%%%


%tabelki
%Create tabular cells spanning multiple rows
%\usepackage{multirow}
\usepackage{multirow}
\usepackage{tabularx}
%array – Extending the array and tabular environments
\usepackage{array}

%%%%%%%%%%%%%%%%
%http://tex.stackexchange.com/questions/191709/biblatex-biber-babel-configuration-for-polish-language-and-style-of-bibliog
\usepackage[
    backend=biber,
    style=authoryear-icomp, %do weryfikacji
    sortlocale=auto, %pl_PL
    natbib=true,
    url=false, % @online is always printed
    doi=true, %The Digital Object Identifier of the work.
    eprint=false, %3.12.7 Electronic Publishing Information
    hyperref=true
]{biblatex}
\addbibresource{biblatex-examples.bib}

%%%%%%%%%%%%%%%






\usepackage{enumitem}	% enumitem.pdf
%\setlist{listparindent=\parindent, parsep=\parskip} % potrzebuje enumitem

\PassOptionsToPackage{hyphens}{url} %break on hyphen http://tex.stackexchange.com/questions/49788/hyperref-url-long-url-with-dashes-wont-break

%%%%%%%%%%%%%%

%advanced hyperref
\usepackage{bookmark}
%package navigator - advaced pdf
\usepackage{listings}
\usepackage[autostyle]{csquotes}
%\DeclareQuoteAlias{croatian}{polish}
%Package csquotes Warning: No style for language 'polish'.
%(csquotes)                Using fallback style on input line 207.
\DeclareQuoteStyle[quotes]{polish}
  {\quotedblbase}
  {\textquotedblright}
  [0.05em]
  {\quotesinglbase}
  {\fixligatures\textquoteright}
\DeclareQuoteAlias[quotes]{polish}{polish}
%\DeclareQuoteOption{polish}
%http://tex.stackexchange.com/questions/191709/biblatex-biber-babel-configuration-for-polish-language-and-style-of-bibliog

%http://tex.stackexchange.com/questions/220395/why-does-hyperref-complain-about-token-not-allowed-in-a-pdf-string-in-my-gloss
\pdfstringdefDisableCommands{\let\textpolish\@firstofone}
\hypersetup{
    pdfencoding=auto,
    pdfauthor={\getautor},
    pdftitle={\gettytul},
    pdfsubject={Praca magisterska},
    pdfkeywords={\getslowakluczowe},
    colorlinks=true,
    %hidelinks
    }
%http://tex.stackexchange.com/questions/18924/pdftex-warning-ext4-destination-with-the-same-identifier-nam-epage-1-has
%%%%%%%%%%%%%%%

%advanced glossaries
%http://tex.stackexchange.com/questions/8946/how-to-combine-acronym-and-glossary

\usepackage[xindy,acronym,toc]{glossaries}
\makeglossaries
\renewcommand*\glspostdescription{\dotfill}

\newacronym[longplural={Frames per Second}]{LB}{LB}{Load-Balancing}
\newacronym[longplural={Frames per Second}]{HA}{HA}{High-Availbility}
\newacronym[longplural={Frames per Second}]{fpsLabel}{FPS}{Frame per Second}

\newacronym{la}{LA}{Los Angeles}
\newacronym{un}{UN}{United Nations}


\newglossaryentry{Linux}
{
  name=Linux,
  description={is a generic term referring to the family of Unix-like
               computer operating systems that use the Linux kernel},
  plural=Linuces
}

\newglossaryentry{SPoF}
{
  name=Single Point of Failure,
  description={is a generic term referring to the family of Unix-like
               computer operating systems that use the Linux kernel},
  plural=Linuces
}

%Don’t use \gls in chapter or section headings as it can have some unpleasant
%side-effects. Instead use \glsentrytext for regular entries and one of
%\glsentryshort, \glsentrylong or \glsentryfull for acronyms



%%%%%%%%%%%%%%%%%%%%%%%%%%%%%%%%%%%%%%%%%%%%%%%%%%%%%%%%%%%%%%%%%%%%%%%%%%%%%
%NOTES
% http://tex.stackexchange.com/questions/9796/how-to-add-todo-notes

\usepackage{xargs}                      % Use more than one optional parameter in a new commands
%\usepackage[dvipsnames]{xcolor}  % Coloured text etc.
% earlier loaded
%\PassOptionsToPackage{dvipsnames}{xcolor}

\usepackage[colorinlistoftodos,prependcaption,textsize=tiny,obeyFinal]{todonotes}
\newcommandx{\unsure}[2][1=]{\todo[linecolor=red,backgroundcolor=red!25,bordercolor=red,#1]{#2}}
\newcommandx{\change}[2][1=]{\todo[linecolor=blue,backgroundcolor=blue!25,bordercolor=blue,#1]{#2}}
\newcommandx{\info}[2][1=]{\todo[linecolor=OliveGreen,backgroundcolor=OliveGreen!25,bordercolor=OliveGreen,#1]{#2}}
\newcommandx{\improvement}[2][1=]{\todo[linecolor=Plum,backgroundcolor=Plum!25,bordercolor=Plum,#1]{#2}}
\newcommandx{\thiswillnotshow}[2][1=]{\todo[disable,#1]{#2}}
% TODO http://tex.stackexchange.com/questions/5180/how-can-i-get-todonotes-in-the-left-margin

%\usepackage{tikz}
\usepackage[printwatermark]{xwatermark}
\newsavebox\mybox
\savebox\mybox{\tikz[color=red,opacity=0.3]\node{draft};}
\newwatermark*[
  allpages,
  angle=45,
  scale=6,
  xpos=-20,
  ypos=15
]{\usebox\mybox}

\usepackage{lipsum}

%\todo[inline]{The original todo note withouth changed colours.\newline Here's another line.}
%\unsure{Is this correct?}\unsure{I'm unsure about also!}
%\change{Change this!}
%\info{This can help me in chapter seven!}
%\improvement{This really needs to be improved!\newline\newline What was I thinking?!}

%\%thiswillnotshow{This is hidden since option `disable' is chosen!}
%\improvement[inline]{The following section needs to be rewritten!}
%\lipsum[11]
%\newpage
%\listoftodos[Notes]


% %%%%%%%%%%%%%%%%%%%%%%%%%%%%%%%%%%%%%%%%%%%%%%%%%%%%%%%%%%%%%%%%%%%%%%%%%%%%
%prevent hyphenation of words on more than two consecutive lines
% avoid hyphenation points that would isolate a single letter
%\usepackage[draft,nosingleletter]{impnattypo}
%\usepackage[draft,nosingleletter,hyphenation,rivers]{impnattypo} %TODO
%wdowy i bękarty
%sprawdzic to defaultlines TODO
\usepackage[defaultlines=4,all]{nowidow}

\pagestyle{fancy}
\begin{document}

% Polecenie \makeinfo produkuje karte informacyjna studenta.
% Sklada sie ja w dziekanacie razem z praca, ale NIE jest
% ona czescia pracy i nie powinna sie w niej znajdowac!
% Po wydrukowanie jednej karty informacyjnej mozna to polecenie
% usunac i drukowac tylko sama prace.
%\makeinfo

% 1 strona tytulowa
\pagenumbering{gobble}
%\thispagestyle{empty}
\maketitle

% 2 opcjonalnie karta pracy dyplomowej, jeżeli wydział stosuje
%\makeinfo

% draft git reflog
\section*{git last changes}
\thispagestyle{empty}
DRAFT: Ta strona będzie drukowane na drugiej karty tytularnej i nie ma na niej numeracji
\lstinputlisting{lastChanges.tex}

% 3 streszczenie pracy
\cleardoublepage
\pagenumbering{roman}
\setcounter{page}{1}
\makeabstracts

% 4 oswiadczenie o autorstwie pracy z USOS
\cleardoublepage
%\thispagestyle{empty}
Oswiadczenie o autorstwie pracy z USOS - do podmiany
% \newpage

% 5 oswiadczenie o udzeleniu licencji - opcjonalne

% 6 spis tresci
\cleardoublepage
\pagestyle{fancy}
\tableofcontents

% 7 kolejne rodzialy
\cleardoublepage
\pagestyle{fancy}
%\setcounter{page}{1}
\pagenumbering{arabic}
% Article top matter
\chapter{Sample chapter}


\section{Introduction}
In \gls{la} this article, I shall discuss some of the fundamental topics in
producing a structured document.  This document itself does not go into
much depth, but is instead the output of an example of how to implement
structure. Its \LaTeX{} source, when in used with my tutorial
(http://www.comp.leeds.ac.uk/andyr/misc/latex/\-latextutorial2.html)
provides all the relevant information.

This small document is designed to illustrate how easy it is to create a
well structured document within \LaTeX\cite{chiu}.  You should quickly be able to
see how the article looks very professional, despite the content being
far from academic.  Titles, section headings, justified text, text
formatting etc., is all there, and you would be surprised when you see
just how little markup was required to get this output.

\section{Structure}
One of the great advantages of \LaTeX{} is that all it needs to know is
the structure of a document, and then it will take care of the layout
and presentation itself.  So, here we shall begin looking at how exactly
you tell \LaTeX{} what it needs to know about your document.

\subsection{Top Matter}
The first thing you normally have is a title of the document, as well as
information about the author and date of publication.  In \LaTeX{} terms,
this is all generally referred to as \emph{top matter}.

Costam \gls{Linux}

\subsubsection{Article Information}
%Set up an 'itemize' environment to start a bulleted list.  Each
%individual item begins with the \item command.  Also note in this list
%that it has two levels, with a list embedded in one of the list items.
\begin{itemize}
	\item \texttt{\textbackslash title\{\emph{title}\}} - The title of the article.
	\item \texttt{\textbackslash date} - The date. Use:
		\begin{itemize}
			\item \texttt{\textbackslash date\{\textbackslash today\}} - to get the
			date that the document is typeset.
			\item \texttt{\textbackslash date\{\emph{date}\}} - for a  %\emph{} emphasises the specified text.  Italics by default.
			specific date.
			\item \texttt{\textbackslash date\{\}} - for no date.
		\end{itemize}
\end{itemize}

\subsubsection{Author Information}
The basic article class only provides the one command:
\begin{itemize}
	\item \texttt{\textbackslash author} - The author of the document.
\end{itemize}
\info[inline]{Notatka}
It is common to not only include the author name, but to insert new
lines (\texttt{\textbackslash\textbackslash}) after and add things such
as address and email details.  For a slightly more logical approach, use
the AMS article class (\texttt{amsart}) and you have the following extra
commands:

\begin{itemize}
	\item \texttt{\textbackslash address} - The author's address.  Use
	the new line command (\texttt{\textbackslash\textbackslash}) for
	line breaks.
	\item \texttt{\textbackslash thanks} - Where you put any acknowledgments.
	\item \texttt{\textbackslash email} - The author's email address.
	\item \texttt{\textbackslash urladdr} - The URL for the author's web page.
\end{itemize}

\subsection{Sectioning Commands}
The commands for inserting sections are fairly intuitive.  Of course,
certain commands are appropriate to different document classes.
For example, a book has chapters but a article doesn't. Odnosnik do tabeli \ref{tab:sampletab}

%A simple table.  The center environment is first set up, otherwise the
%table is left aligned.  The tabular environment is what tells Latex
%that the data within is data for the table.
\begin{table}
  \caption{tytul tabeli \thetable}
  \label{tab:sampletab}
\begin{center}
	\begin{tabular}{| l | l |}
	%The tabular environment is what tells Latex that the data within is
	%data for the table.  The arguments say that there will be two
	%columns, both left justified (indicated by the 'l', you could also
	%have 'c' or 'r'.  The bars '|' indicate vertical lines throughout
	%the table.

		\hline  % Print horizontal line
		Command & Level \\ \hline  % Columns are delimited by '&'.  And
                                   %rows are delimited by '\\'
		\texttt{\textbackslash part\{\emph{part}\}} & -1 \\
		\texttt{\textbackslash chapter\{\emph{chapter}\}} & 0 \\
		\texttt{\textbackslash section\{\emph{section}\}} & 1 \\
		\texttt{\textbackslash subsection\{\emph{subsection}\}} & 2 \\
		\texttt{\textbackslash subsubsection\{\emph{subsubsection}\}} & 3 \\
		\texttt{\textbackslash paragraph\{\emph{paragraph}\}} & 4 \\
		\texttt{\textbackslash subparagraph\{\emph{subparagraph}\}} & 5 \\
		\hline
	\end{tabular}
\end{center}
  \caption*{zrodlo tabeli}
\end{table}

\unsure[inline]{Fix me}

Numbering of the sections is performed automatically by \LaTeX{}, so don't
bother adding them explicitly, just insert the heading you want between
the curly braces.  If you don't want sections number, then add an asterisk (*) after the
section command, but before the first curly brace, e.g., \texttt{\textbackslash
section*\{A Title Without Numbers\}}.

\section{Image section}

\begin{figure}[h]
\begin{center}
    \includegraphics[width=0.5\textwidth]{foto.png}
\end{center}
  \caption[short - gull]{A picture of a gull. bardzo dlugi opis rysnklu jeszcze dluzej aaaaaaaaaaaaaaaaaa rysunkek bla bla bla bla blaadasd sadnasujvndsjiogf jdifjsdiof dsjfiods jfido \thefigure}
  \caption*{zrodlo rysunku \ref{ap:ap1}}
  \label{img:gull}
\end{figure}

\chapter{chapter2}

\section{Image section2}

\begin{figure}[h]
\begin{center}
    \includegraphics[width=0.5\textwidth]{foto.png}
\end{center}
  \caption[short - gull2]{A picture of a gull. bardzo dlugi opis rysnklu jeszcze dluzej aaaaaaaaaaaaaaaaaa rysunkek bla bla bla bla bla asdkaosdkaosdkao dsakods akdosa dkaos dkaso \thefigure}
  \caption*{zrodlo rysunku \ref{ap:ap1}}
  \label{img:gull2}
\end{figure}

\section{section name}
\begin{itemize}
\color{blue}
\item First item
\item Second item
\end{itemize}



%Create the environment for the bibliography.  Since there is only one
%reference, set the label width to be one character (I shall follow
%convention as use the number '9'.  This is because it helps to remind
%that it is the maximum number of refs that is now permitted by that
%width).




% 8 bibliografia
%\bibliographystyle{alpha}
%\bibliographystyle{plain}
%\bibliography{mgr}
\cleardoublepage
%\setcounter{page}{1}
%\pagenumbering{roman}
\label{bibliography}
\printbibliography[heading=bibintoc]

% 9 wykaz symboli i skrotow
\clearpage
%\printglossary
%\printglossaries
%\renewcommand{\glossarypreamble}{}
\printglossary[type=\acronymtype,title=Wykaz skrótów]
\printglossary[type=main]

% 10 spis rysunkow
\clearpage\phantomsection
\label{listoffigures}
\listoffigures
\addcontentsline{toc}{chapter}{\listfigurename}

% 11 spis tabel
\clearpage\phantomsection
\label{listoftables}
\listoftables
\addcontentsline{toc}{chapter}{\listtablename}


% 12 spis zalacznikow
% 13 zalaczniki
\clearpage\phantomsection
%logic
%\thispagestyle{empty} %plain ?
\newcommand{\listappendicesname}{Spis załączników}
\newlistof{appendices}{apc}{\listappendicesname}
\newcommand{\appendices}[1]{\addcontentsline{apc}{appendices}{#1}}
%\newcommand{\newappendix}[1]{\subsection*{#1}\appendices{#1}}
\newcommand{\newappendix}[1]{\section{#1}\appendices{#1}}
%\addcontentsline{toc}{chapter}{\listappendicesname}
\renewcommand{\cftapctitlefont}{\bfseries\LARGE}

%print list
\listofappendices
%schema for enumerating
\label{listof}
%\chapter*{\listappendicesname}



%add app
\appendix
\addtocontents{toc}{\protect\setcounter{tocdepth}{0}}

\chapter*{Załączniki}
\renewcommand{\thesection}{Załącznik \Alph{section}:}

\newappendix{ap1}
\label{ap:ap1}
asdasd

\addcontentsline{toc}{chapter}{\listappendicesname}

%\addtocontents{toc}{\protect\contentsline{chapter}{\listfigurename}{\getpagerefnumber{listoffigures}}{\indexname}}
%\addtocontents{toc}{\protect\contentsline{chapter}{\listfigurename}{\getpagerefnumber{bibliography}}{}}
%\addtocontents{toc}{\protect\contentsline{chapter}{\listtablename}{\getpagerefnumber{listoftables}}{}}
%\addtocontents{toc}{\protect\contentsline{chapter}{\listappendicesname}{\getpagerefnumber{listofappendices}}{}}
%http://tex.stackexchange.com/questions/321132/add-list-of-figures-tables-etc-on-different-position-in-table-of-contents
%http://tex.stackexchange.com/questions/194720/hyperref-in-toc-without-page-number


\end{document}
